\usepackage{microtype}
\usepackage{stmaryrd}
\usepackage{float}
\usepackage{adjustbox}
\usepackage{amsfonts}
\usepackage{amssymb}
\usepackage{amsmath,bbm}
\usepackage{xspace}
\usepackage[normalem]{ulem}
\usepackage{hyperref}
\usepackage[capitalise]{cleveref}
\usepackage{graphicx}
\usepackage{wasysym}
\usepackage{mathtools}
\usepackage{array,amsmath}
% Editorial wrapper: transparent omitted; no author transparent command; TikZ fadings provides pgftransparent.
\usepackage{marvosym}
\usepackage{upgreek}
\DeclareMathAlphabet\mathbfcal{OMS}{cmsy}{b}{n}

\usetikzlibrary{trees}
\usetikzlibrary{shapes,fit,backgrounds,fadings}
\usetikzlibrary{positioning,arrows,calc,decorations.markings}
\tikzset{
modal/.style={>=stealth',shorten >=1pt,shorten <=1pt,auto,node distance=1.5cm,semithick},
world/.style={circle,draw,minimum size=0.5cm,fill=gray!15},
point/.style={circle,draw,inner sep=0.5mm,fill=black},
reflexive above/.style={->,loop,looseness=7,in=120,out=60},
reflexive below/.style={->,loop,looseness=7,in=240,out=300},
reflexive left/.style={->,loop,looseness=7,in=150,out=210},
reflexive right/.style={->,loop,looseness=7,in=30,out=330}
}
\newcommand{\hyperedgei}[4][180]{
     \draw (#2.#1) ++(#1:.5) node[] {$\pred{R}$} edge (#2) edge[->] (#3) edge[->] (#4);    
}
\newcommand{\hyperedge}[4][180]{
     \draw (#2.#1) ++(#1:.5) node[] {} edge (#2) edge[->] (#3) edge[->] (#4);    
}

\makeatletter
\pgfdeclareverticalshading{pgf@lib@fade@South}{100bp}
{color(0bp)=(pgftransparent!100); color(40bp)=(pgftransparent!100);
 color(75bp)=(pgftransparent!0); color(100bp)=(pgftransparent!0)}%
\pgfdeclarefading{South}{\pgfuseshading{pgf@lib@fade@South}}%
\makeatother

%\usepackage{balance} % balance columns on last page
\usepackage{bold-extra} % for bold typewriter fonts
% Editorial wrapper: unavailable style-only nimbusmononarrow omitted; standard author mathtt remains.
\usepackage[scr=boondoxo]{mathalpha}
\usepackage{scalerel} % allows to scale symbols 

\hyphenation{tree-width clique-width ho-mo-mor-phism iso-mor-phism}

% Environments
%% I changed the style for the theorem environments to the lmcs style. If you do not agree with this, let us know
\theoremstyle{plain}
\newaliascnt{theorem}{thm}
\newtheorem{theorem}[theorem]{Theorem}
\aliascntresetthe{theorem}
\newaliascnt{lemma}{thm}
\newtheorem{lemma}[lemma]{Lemma}
\aliascntresetthe{lemma}
\newaliascnt{shorttheorem}{thm}
\newtheorem{shorttheorem}[shorttheorem]{Thm.}
\aliascntresetthe{shorttheorem} % emergency use only 
\newaliascnt{proposition}{thm}
\newtheorem{proposition}[proposition]{Proposition}
\aliascntresetthe{proposition}
\newaliascnt{corollary}{thm}
\newtheorem{corollary}[corollary]{Corollary}
\aliascntresetthe{corollary}
\newaliascnt{observation}{thm}
\newtheorem{observation}[observation]{Observation}
\aliascntresetthe{observation}

\theoremstyle{definition}
\newaliascnt{definition}{thm}
\newtheorem{definition}[definition]{Definition}
\aliascntresetthe{definition}
\newaliascnt{assumption}{thm}
\newtheorem{assumption}[assumption]{Assumption}
\aliascntresetthe{assumption}
\newaliascnt{example}{thm}
\newtheorem{example}[example]{Example}
\aliascntresetthe{example}
\newaliascnt{remark}{thm}
\newtheorem{remark}[remark]{Remark}
\aliascntresetthe{remark}
\newaliascnt{folk}{thm}
\newtheorem{folk}[folk]{Fact}
\aliascntresetthe{folk}
\newaliascnt{property}{thm}
\newtheorem{property}[property]{Property}
\aliascntresetthe{property}


\newenvironment{subproof}[1][\proofname]{%
  \renewcommand{\qedsymbol}{$\blacksquare$}%
  \begin{proof}[#1]%
}{%
  \end{proof}%
}

%\newcommand{\remarkstart}{\color{blue}}
\newcommand{\remarkstart}{\color{black}}


\newtheoremstyle{myexample}% ⟨name⟩ 
{3pt}% ⟨Space above⟩ 
{3pt}% ⟨Space below⟩
{\rm}% ⟨Body font⟩
{}% ⟨Indent amount⟩
{\bf}% ⟨Theorem head font⟩
{.}% ⟨Punctuation after theorem head⟩
{.5em}% ⟨Space after theorem head⟩
{}% ⟨Theorem head spec (can be left empty, meaning ‘normal’)⟩
\theoremstyle{myexample}
\newaliascnt{xample}{theorem}
\newtheorem{xample}[xample]{Example}
\aliascntresetthe{xample}
\newaliascnt{emark}{theorem}
\newtheorem{emark}[emark]{Remark}
\aliascntresetthe{emark}


\usepackage{framed}
\definecolor{shadecolor}{rgb}{0.93,0.93,0.93}
\definecolor{darkerblue}{rgb}{0.0,0.0,0.6}
\definecolor{darkblue}{rgb}{0.0,0.0,0.9}
\definecolor{darkerred}{rgb}{0.6,0.0,0.0}
\definecolor{darkred}{rgb}{0.9,0.0,0.0}
\definecolor{darkergreen}{rgb}{0.0,0.4,0.0}
\definecolor{darkgreen}{rgb}{0.0,0.6,0.0}
\definecolor{darkercyan}{rgb}{0.0,0.4,0.6}
\definecolor{darkcyan}{rgb}{0.0,0.6,0.9}
\renewenvironment{example}%[1][]
{
	\begin{shaded*}	
		\vspace{-1.5ex}
		\begin{exa}%[#1]
		}{
		\end{exa}
		\vspace{-1.5ex}
	\end{shaded*}
}
\renewenvironment{remark}%[1][]
{
	\begin{shaded*}	
		\vspace{-1.5ex}
		\begin{rem}%[#1]
		}{
		\end{rem}
		\vspace{-1.5ex}
	\end{shaded*}
}

\crefname{rem}{Remark}{Remarks}
\crefname{exa}{Example}{Examples}

\newcommand{\footnotetextbefore}[1]{\addtocounter{footnote}{1}\footnotetext{#1}\addtocounter{footnote}{-1}
}

%%%%%%%%%%%%%%%%%%%%%%%%%%%%%%%%%%%%%%%%%%%%%%%%%%%%%%%%%
%%%%%%%%%%%%%%%%%%%NOTATION SECTION%%%%%%%%%%%%%%%%%%%%%%
